% ===========================================================================
\section{Decomposing the student's gap}\label{sec:decomposition}
% ===========================================================================

\subsection{Pipeline and master decomposition}

\paragraph{Pipeline.}
We draw $n$ prompts $x_i\sim\cD$ i.i.d. For each prompt the teacher generates $y^{(i)}\sim P^{x_i}$ and
supplies either soft labels $p_t(\cdot\mid s^{(i)}_t)$ at every position, or only the sampled tokens (hard
labels). On-policy variants additionally sample $y\sim Q^{x_i}_\theta$. Training uses a prefix law
$\nu(\theta)=(\nu_t(\theta))_t$: $\nu_t=\dteach t$ (off-policy), $\nu_t=\dstud t$ (on-policy), or a mixture
$\nu_t=\lambda\dstud t+(1-\lambda)\dteach t$ \citep{agarwal2024gkd}. The per-state loss is
$\ell_\theta(s)=\mathsf D(p_t(\cdot\mid s)\|q_t(\cdot\mid s))$. The population and empirical training
objectives are
\[
F(\theta):=\E_{x\sim\cD}\sum_{t=1}^T\E_{s\sim\nu_t(\theta)\mid x}\ell_\theta(s),\qquad
\hat F_n(\theta):=\frac1n\sum_{i=1}^n\sum_{t=1}^T\ell_\theta(s^{(i)}_t).
\]
For on-policy terms, $\hat F_n$ is defined with fixed sampling noise (for example Gumbel-max
reparametrisation), so that it is a fixed random function of $\theta$. Let
$\theta^\star\in\argmin_{\Theta_\TS}F$, $\htheta\in\argmin\hat F_n$, and let $\ttheta$ be the output of the
optimiser. The \emph{evaluation functional} $\cE_\mu(\theta)$ is a sequence-level criterion under the
deployment law, for example $\E_\mu\KL(\Pseq\|\Qseq)$ or a task regret. $\cE_\cD$ is the same functional
under $\cD$.

\begin{theorem}[Master decomposition]\label{thm:master}
\Pv
\begin{equation}\label{eq:master}
\cE_\mu(\ttheta)=\underbrace{F(\theta^\star)}_{\cA\ \text{approximation}}
+\underbrace{F(\ttheta)-F(\theta^\star)}_{\text{estimation}+\text{optimisation}}
+\underbrace{\cE_\cD(\ttheta)-F(\ttheta)}_{\cE_{\mathrm{shift}}\ \text{prefix shift}}
+\underbrace{\cE_\mu(\ttheta)-\cE_\cD(\ttheta)}_{\cE_{\mathrm{cov}}\ \text{covariate shift}},
\end{equation}
\begin{equation}\label{eq:est-opt}
F(\ttheta)-F(\theta^\star)\le\underbrace{2\sup_{\theta}|\hat F_n(\theta)-F(\theta)|}_{\cE_{\mathrm{est}}}
+\underbrace{\hat F_n(\ttheta)-\hat F_n(\htheta)}_{\cE_{\mathrm{opt}}}.
\end{equation}
\end{theorem}
\begin{proof}
\eqref{eq:master} telescopes. For \eqref{eq:est-opt}, write
$F(\ttheta)-F(\theta^\star)=[F-\hat F_n](\ttheta)+[\hat F_n(\ttheta)-\hat F_n(\htheta)]+[\hat F_n(\htheta)-\hat F_n(\theta^\star)]+[\hat F_n-F](\theta^\star)$,
and note that the third bracket is $\le0$. This is the approximation--estimation--optimisation split of
\citet{bottou2008tradeoffs} with two shift terms added.
\end{proof}

\noindent
When $\cE$ is a KL, the behavioural criterion follows from Pinsker:
$\E_\mu\TV(\Pseq,Q^x_{\ttheta})\le\sqrt{\cE_\mu(\ttheta)/2}$.
The key structural fact is that $\cE_{\mathrm{shift}}\equiv0$ for exactly two pairings (\cref{thm:chain}). The
table lists what each pairing's training objective $F$ equals and the resulting bound on sequence-level
behaviour:
\begin{center}\small
\begin{tabular}{@{}p{0.17\linewidth}p{0.12\linewidth}p{0.3\linewidth}p{0.3\linewidth}@{}}
\toprule
per-token loss & prefix law & $F$ equals & $\E_\cD\TV(P,Q_\theta)\le$\\\midrule
forward KL & teacher & $\E_\cD\KL(P\|Q_\theta)$; $\cE_{\mathrm{shift}}=0$ & $\sqrt{F/2}$\\
reverse KL & student & $\E_\cD\KL(Q_\theta\|P)$; $\cE_{\mathrm{shift}}=0$ & $\sqrt{F/2}$\\
forward KL & student & hybrid (\cref{prop:hybrid-div}) & $\sqrt{TF/2}$\\
reverse KL & teacher & hybrid (\cref{prop:hybrid-div}) & $\sqrt{TF/2}$\\
TV & either & hybrid (\cref{prop:hybrid}) & $F$\\
$\JSD_\beta$ & either & hybrid & $\sqrt{TF/(2\beta(1-\beta))}$\\\bottomrule
\end{tabular}
\end{center}
For the hybrid rows, the factor $T$ inside the square root comes from our Cauchy--Schwarz argument. A
sequential subadditivity inequality for squared Hellinger distance \citep{foster2021statistical}
improves it to $O(\log T)$; we have not re-derived that lemma's constant. For hybrid rows, the on-policy
per-token loss can also exceed the training loss. For bounded losses the excess is at most
$G\sum_t\sum_{k<t}e^{\mathrm{off},\TV}_k$, which is $O(T^2)$ in the per-token TV error (\cref{prop:exposure}). For reverse KL
on teacher prefixes it is unbounded relative to the training loss.

\subsection{Approximation error}\label{sec:approx}

\begin{lemma}[Log-ratio variance]\label{lem:logratio}
\Pv\Ck Let $\delta=\log(p/q)$ on $\operatorname{supp}p$ and $M:=\max\{0,\max_y\delta(y)\}$. Then
$\KL(p\|q)\ge\frac{e^{-M}}2\E_p[\delta^2]\ge\frac{e^{-M}}2\Var_p(\delta)$.
\end{lemma}
\begin{proof}
$\E_p[e^{-\delta}]=q(\operatorname{supp}p)\le1$. Taylor expansion gives $e^{-\delta}=1-\delta+\frac{\delta^2}2e^{-\xi}$
with $\xi$ between $0$ and $\delta$, so $e^{-\xi}\ge e^{-M}$. Taking $\E_p$:
$1\ge1-\KL(p\|q)+\frac{e^{-M}}2\E_p\delta^2$.
\end{proof}

\begin{theorem}[Softmax-bottleneck approximation floor]\label{thm:bottleneck}
\Pa{\cref{as:linear}}\Ck Let $s_1,\dots,s_N$ be contexts at which the teacher has full support. Let
$A\in\R^{N\times V}$ have rows $\log p(\cdot\mid s_i)$, let $\pi_i:=\min_yp(y\mid s_i)$, and let
$\Pi=\diag(\pi_i)$. For every student of width $d_\TS$ with $\max_{i,y}\log\frac{p(y\mid s_i)}{q_\theta(y\mid s_i)}\le M$,
\[
\frac1N\sum_{i=1}^N\KL\big(p(\cdot\mid s_i)\,\|\,q_\theta(\cdot\mid s_i)\big)\;\ge\;\frac{e^{-M}}{2N}\sum_{j>d_\TS+1}\sigma_j^2\big(\Pi^{1/2}AC\big).
\]
\end{theorem}
\begin{proof}
By \cref{lem:logratio}, $\KL_i\ge\frac{e^{-M}}2\Var_{p_i}(\delta_i)$. Also
$\Var_{p_i}(\delta_i)=\min_c\E_{p_i}(\delta_i-c)^2\ge\pi_i\min_c\|\delta_i-c\one\|^2=\pi_i\|C\delta_i\|^2$.
The student's log-probability rows are $Wh_i+b-\mathrm{lse}_i\one$, so the centred matrix $ZC$ has rows
$C(Wh_i+b)$ and rank at most $d_\TS+1$. Left-multiplication by the diagonal $\Pi^{1/2}$ does not
increase rank. Hence
$\sum_i\pi_i\|C\delta_i\|^2=\|\Pi^{1/2}(A-Z)C\|_F^2\ge\min_{\rank(Y)\le d_\TS+1}\|\Pi^{1/2}AC-Y\|_F^2$, which
equals the tail sum by Eckart--Young--Mirsky \citep{eckart1936approximation}.
\end{proof}

\begin{remark}\label{rem:bottleneck}
(i)~If the teacher also has a linear read-out of width $d_\TT$, then $\rank(AC)\le d_\TT+1$. The floor is
then a sum over the teacher's singular values with indices in $(d_\TS+1,d_\TT+1]$, and it vanishes if
$d_\TS\ge d_\TT$. This quantifies the ``softmax bottleneck'' \citep{yang2018breaking}. Its relevance to
small language models was studied empirically by \citet{godey2024small}; we are fairly, but not fully,
confident of that reference's exact claims.
(ii)~The factor $\pi_i$ is crude. Replacing $\Pi^{1/2}C$ by any common lower bound $\Gamma^{1/2}$ on the
Fisher matrices $\diag(p_i)-p_ip_i^\top$ gives the tail spectrum of $A\Gamma^{1/2}$. In the check of
\texttt{verification/check\_capacity.py}, the row-weighted floor is within a factor $15$ of a fitted
student's loss.
(iii)~The floor is a statement about the \emph{read-out}. When $d_\TS\ge d_\TT$ it is silent;
information-theoretic floors (\cref{sec:lower}) take over.
\end{remark}

\begin{corollary}[Approximation grows linearly with the horizon]\label{cor:approx-T}
\Pv If the per-token approximation floor on teacher prefixes is $\ge\eps_0$ at every $t$, then
$\min_\theta\E_\cD\KL(P\|Q_\theta)\ge T\eps_0$. Likelihood losslessness at level $\eps$ therefore requires
$\eps_0\le\eps/T$. For TV, \cref{prop:hybrid} gives only $\E\TV(P,Q_\theta)\ge\frac12\max_t\E_{\dteach t}\TV_t$.
Whether TV grows with $T$ depends on how the errors are structured (\cref{ex:regimes}).
\end{corollary}

\subsection{Estimation error}\label{sec:estimation}

\begin{theorem}[Uniform deviation, soft-label off-policy KD]\label{thm:unif}
\Pa{\cref{as:bounded}} With soft labels, $\ell(x,y;\theta)=\sum_t\KL(p_t(s_t)\|q_t(s_t))\in[0,TB]$, since
$\KL(p\|q)\le\E_p[-\log q]\le B$. With probability $\ge1-\delta$,
\[
\sup_\theta|\hat F_n(\theta)-F(\theta)|\le2\,\Rad_n(\ell\circ\Theta_\TS)+TB\sqrt{\tfrac{\log(2/\delta)}{2n}}.
\]
For a finite class (\cref{as:finite}), $\Rad_n\le TB\sqrt{2R\ln2/n}$ by Massart's lemma.
\end{theorem}
\begin{proof}
Symmetrisation and McDiarmid \citep{bartlett2002rademacher}.
\end{proof}
\noindent
This bound is pessimistic: a slow rate, and linear in $T$. For the realisable case, a
horizon-free fast rate holds even with \emph{hard} labels. The logic is due to
\citet{zhang2006epsilon} and \citet{wong1995probability}; its horizon-independence in imitation learning
was emphasised by \citet{foster2024behavior}.

\begin{theorem}[Horizon-free rate for sequence-level hard-label KD]\label{thm:mle}
\Pa{\cref{as:real,as:finite}} Let $\htheta$ maximise $\sum_{i=1}^n\log Q_\theta(y^{(i)}\mid x_i)$ over
$\Theta_\TS$, with $y^{(i)}\sim P^{x_i}$. With probability $\ge1-\delta$,
\begin{gather*}
\E_{x\sim\cD}\Hel^2\big(\Pseq,Q^x_{\htheta}\big)\le\frac{\log(|\Theta_\TS|/\delta)}n\le\frac{R\ln2+\log(1/\delta)}n,\\
\E_{x\sim\cD}\TV\big(\Pseq,Q^x_{\htheta}\big)\le\sqrt{\frac{2(R\ln2+\log(1/\delta))}n}.
\end{gather*}
Neither bound depends on $T$.
\end{theorem}
\noindent\emph{Proof.} See \cref{app:mle}. The key lemma is
$\E_{y\sim P}\sqrt{Q(y)/P(y)}=1-\Hel^2(P,Q)\le e^{-\Hel^2}$, combined with a Chernoff bound and a union
bound over $\Theta_\TS$.

\begin{proposition}[Soft labels Rao--Blackwellise hard labels]\label{prop:rb}
\Pv Fix a state law $\nu$ with states drawn independently of labels. Let $\hat h=-\log q_\theta(y\mid s)$ and
$\hat s=\mathrm{CE}(p(\cdot\mid s),q_\theta(\cdot\mid s))$, where $s\sim\nu$ and $y\sim p(\cdot\mid s)$. Then
$\E\hat h=\E\hat s$ and
\begin{gather*}
\Var\hat h-\Var\hat s=\E_{s\sim\nu}\Var_{y\sim p(\cdot\mid s)}\log q_\theta(y\mid s)\ \ge0,\\
\E\|\nabla\hat h-\nabla\E\hat h\|^2-\E\|\nabla\hat s-\nabla\E\hat s\|^2=\E_\nu\tr\Cov_{y\sim p(\cdot\mid s)}\nabla_\theta\log q_\theta(y\mid s).
\end{gather*}
At $\theta=\theta^\circ$ (\cref{as:real}), the gradient gap equals $\E_\nu\tr\mathcal I(s)$. Here
$\mathcal I(s)=J(s)^\top(\diag p-pp^\top)J(s)$ is the per-state Fisher information and $J=\partial z/\partial\theta$.
The gap vanishes wherever $p(\cdot\mid s)$ is a point mass.
\end{proposition}
\begin{proof}
$\hat s=\E[\hat h\mid s]$, so the law of total variance applies; likewise for the gradient, since
$\nabla\hat s=\E[\nabla\hat h\mid s]$. At $\theta^\circ$,
$\Cov_{y\sim p}\nabla\log q=J^\top\Cov_{y\sim p}(e_y)J=J^\top(\diag p-pp^\top)J$.
\end{proof}
\begin{remark}
For whole teacher-sampled sequences, later states depend on earlier labels. Position-wise
Rao--Blackwellisation still holds, but cross-time covariances can in principle make the variance of the
\emph{sum} larger. We therefore claim no sequence-level variance ordering. The variance-reduction view of
distillation is developed by \citet{menon2021statistical}. Related analyses include
\citet{lopezpaz2016unifying}, \citet{phuong2019towards}, \citet{dao2021knowledge} and
\citet{harutyunyan2023supervision}.
\end{remark}

\begin{proposition}[Soft labels: exact recovery versus $k/2n$]\label{prop:exact}
Let $q_\Theta(\cdot\mid s)=\softmax(\Theta\phi(s))$ with $\Theta\in\R^{V\times d}$, and let $p=q_{\Theta^\circ}$.
\begin{enumerate}[label=(\alph*),nosep]
\item \Pv If $\phi(s_1),\dots,\phi(s_n)$ span $\R^d$ (so $n\ge d$), every minimiser of
$\sum_i\KL(p(\cdot\mid s_i)\|q_\Theta(\cdot\mid s_i))$ satisfies $q_\Theta=p$ at \emph{every} $s$.
Estimation error is zero.
\item \Kn If $\E[\phi\phi^\top]\succ0$, the hard-label MLE from $n$ i.i.d.\ pairs $(s_i,y_i)$ satisfies
$n\,\KL(p\|q_{\hat\Theta})\Rightarrow\frac12\chi^2_{(V-1)d}$ in distribution. The limiting mean is $\frac{(V-1)d}2$. This is the
classical $k/2n$ law for a regular, well-specified model with $k$ free parameters, the computation behind
AIC \citep{akaike1974new}; see \citet{white1982maximum} for the misspecified case. The \emph{expectation}
$\E\KL$ is infinite for every $n$, because with positive probability the MLE does not exist or puts zero
mass on an observed class. The $k/2n$ law is therefore a statement about the limit law, not about
$\E\KL$.
\end{enumerate}
\end{proposition}
\begin{proof}[Proof of (a)]
The objective is $0$ at $\Theta^\circ$. It vanishes at $\Theta$ iff $(\Theta-\Theta^\circ)\phi(s_i)\in\operatorname{span}(\one)$
for all $i$, which means $C(\Theta-\Theta^\circ)\phi(s_i)=0$. By spanning, $C(\Theta-\Theta^\circ)=0$. So the
logits differ by a constant vector at every $s$, and the softmax is unchanged.
\end{proof}
\noindent
In the realisable case, a soft label is therefore worth ``infinitely many'' hard labels at its state.
Once the teacher's labels are exact, what remains is extrapolation, which is controlled by coverage of
the feature space and by approximation, not by label noise.

\subsection{Optimisation error}\label{sec:optimisation}

\begin{proposition}[Curvature in logit space]\label{prop:convexity}
\Pv\Ck
\begin{enumerate}[label=(\roman*),nosep]
\item $z\mapsto\KL(p\|\softmax z)$ is convex. Its Hessian satisfies
$0\preceq\diag(q)-qq^\top\preceq\frac12C$ \citep{bohning1992multinomial}, so it is $\frac12$-smooth.
\item $z\mapsto\KL(\softmax z\|p)$ is not convex, for any $V\ge2$ and any $p$ of full support. With $V=2$,
$f(z)=\KL(\sigma(z)\|\mathrm{Bern}(p))$ has $f''(z)=\sigma'(z)\,[1+(1-2\sigma(z))(z-z_p)]$ with $z_p=\operatorname{logit}p$, which is
negative whenever $(2\sigma(z)-1)(z-z_p)>1$; for example $p=0.3$, $z=z_p+3$. For general $V$, along the ray
$z=(t,0,\dots,0)$ the second derivative is $-(V-1)(t-2)e^{-t}+O(e^{-t})<0$ for large $t$.
\item If $p$ has full support, then for $\beta$ close to $1$, $z\mapsto\JSD_\beta(p\|\softmax z)$ is not convex:
$\JSD_\beta/(1-\beta)\to\KL(q\|p)$ together with its second derivatives, uniformly on compacts.
\item Even with logits linear in a scalar $\theta$, reverse KL can have spurious local minima. Take two
contexts, $V=2$, logits $z_1=5.2-\theta/2$ and $z_2=-1.5-\theta/3$, and teacher probabilities $0.79$ and $0.26$ for the
first token. The summed reverse KL has strict local minima at $\theta\approx-1.00$ (value $0.2191$) and
$\theta\approx7.23$ (value $0.2300$), separated by a barrier of height $0.2678$ at $\theta\approx3.61$. The forward KL has a
single minimum.
\end{enumerate}
Consequently, for any student whose logits are linear in $\theta$ (for example, training only the
read-out of \cref{as:linear}), off-policy forward-KL distillation is a convex problem. Reverse-KL
objectives can have spurious local minima even then (iv).
\end{proposition}
\begin{proof}
(i) The Hessian of $\mathrm{lse}$ is $\diag(q)-qq^\top$; Böhning's bound gives the smoothness constant.
(ii) $f'(z)=\sigma(z)(1-\sigma(z))(z-z_p)$, because $\log\frac{q}{1-q}=z$; differentiate once more. On the ray,
$q_j=e^{-t}(1+O(e^{-t}))$ for $j>1$, and the reverse KL equals $\log\frac1{p_1}-(V-1)te^{-t}+O(e^{-t})$; differentiate twice.
(iii) Write $\JSD_\beta/(1-\beta)=\KL(q\|M_\beta)+\frac{\beta}{1-\beta}\KL(p\|M_\beta)$. Since $p>0$, $M_\beta\ge\beta p$ is bounded away
from $0$, so $\KL(q\|M_\beta)$ is jointly smooth in $(\beta,z)$ up to $\beta=1$, where it equals $\KL(q\|p)$. Let
$\epsilon=1-\beta$ and $g(s):=\KL(p\|p+s(q-p))$. Then $g(0)=g'(0)=0$, so $g(\epsilon)/\epsilon=\epsilon\int_0^1(1-t)\,g''(\epsilon t)\,dt$. This is
jointly smooth in $(\epsilon,z)$ and vanishes at $\epsilon=0$ together with its $z$-derivatives. Hence second
$z$-derivatives converge uniformly on compacts, and the strict negative curvature of (ii) persists for
$\beta$ near $1$. (iv) is a direct computation
(\texttt{check\_divergences.py}).
\end{proof}

\begin{theorem}[Stochastic optimisation]\label{thm:sgd}
\Kn Under \cref{as:opt}, SGD with a tuned constant step size satisfies
$\min_{k<K}\E\|\nabla F(\theta_k)\|^2\lesssim\frac{L\Delta}K+\sigma\sqrt{\frac{L\Delta}K}$ with
$\Delta=F(\theta_0)-\inf F$ \citep{ghadimi2013stochastic}. If $F$ is convex (\cref{prop:convexity}), then
$\E F(\bar\theta_K)-\min F\lesssim\frac{L\|\theta_0-\theta^\star\|^2}K+\frac{\sigma\|\theta_0-\theta^\star\|}{\sqrt K}$.
In both cases $\sigma^2$ enters, so the Rao--Blackwell gap of \cref{prop:rb} translates directly into
lower optimisation error at fixed compute.
\end{theorem}
\begin{remark}
No non-convex guarantee beyond stationarity is available for deep students. Empirically, optimisation
is often the binding constraint on fidelity. \citet{stanton2021does} find that students frequently fail to
match their teachers even when they have the capacity to do so. \citet{beyer2022knowledge} find that long,
``patient and consistent'' schedules are needed. We treat $\cE_{\mathrm{opt}}$ as measurable
(\cref{sec:experiments}), not as provably small.
\end{remark}

\subsection{Distribution shift and exposure bias}\label{sec:shift}

Let $c_t:\cS_t\times\cV\to[0,1]$ be per-step costs and $J_c(\pi):=\E_{x\sim\cD}\E_{y\sim\pi}\sum_{t=1}^Tc_t(s_t,y_t)$.
Let $r:\cX\times\cY\to[0,1]$ be a terminal reward and $J_r(\pi):=\E\,r(x,y)$. The teacher's cost-to-go is
$Q^p_t(s,a):=c_t(s,a)+\E_p[\sum_{k>t}c_k\mid s_t=s,y_t=a]$, and its \emph{recoverability constant} is
\[
u:=\max_t\sup_s\big[\max_aQ^p_t(s,a)-\min_aQ^p_t(s,a)\big]\in[0,T].
\]

\begin{theorem}[Terminal rewards: linear in $T$]\label{thm:terminal}
\Pv $|J_r(p)-J_r(q_\theta)|\le\E_\cD\TV(P,Q_\theta)\le T\min\{\eoff^{\TV},\eon^{\TV}\}$. Also
$|J_r(p)-J_r(q_\theta)|\le\sqrt{T\eoff^{\KL}/2}$.
\end{theorem}
\begin{proof}
\cref{prop:tv-task,prop:hybrid,cor:seq-tv}.
\end{proof}

\begin{theorem}[Cumulative costs, off-policy: quadratic, and tight]\label{thm:quadratic}
\Pv\Ck
\[
|J_c(p)-J_c(q_\theta)|\le\sum_{t=1}^T(T-t+1)\,e^{\mathrm{off},\TV}_t\le T^2\eoff^{\TV}.
\]
The order is attained. Take a deterministic teacher emitting $a^\star$ in every ``on-track'' state. Let
every deviation lead to absorbing off-track states that cost $1$ per step. Let the student deviate with
probability $\eps$ on track. Then $\eoff^{\TV}=\eps$ and
$J_c(q_\theta)-J_c(p)=\sum_{t=1}^T[1-(1-\eps)^t]=\eps\frac{T(T+1)}2(1+O(\eps T))$.
\end{theorem}
\begin{proof}
$|\E_pc_t-\E_{q}c_t|\le\TV$ between the laws of $y_{\le t}$, which is at most
$\sum_{k\le t}e^{\mathrm{off},\TV}_k$ by the hybrid argument (\cref{prop:hybrid}) applied to the first $t$
tokens. Sum over $t$. In the example, the student is on track at step $t$ with probability $(1-\eps)^{t-1}$.
\end{proof}

\begin{theorem}[On-policy: linear with recoverability]\label{thm:pdl}
\Pv\Ck The performance-difference identity \citep{kakade2002approximately} holds:
\[
J_c(q_\theta)-J_c(p)=\sum_{t=1}^T\E_{s\sim\dstud t}\Big[\E_{a\sim q_t(\cdot\mid s)}Q^p_t(s,a)-V^p_t(s)\Big],\qquad V^p_t(s):=\E_{a\sim p_t(\cdot\mid s)}Q^p_t(s,a).
\]
Hence $|J_c(q_\theta)-J_c(p)|\le u\,T\,\eon^{\TV}$ and, by Pinsker per token and Jensen,
$|J_c(q_\theta)-J_c(p)|\le u\,T\sqrt{\eon^{\mathrm{rKL}}/2}$.
\end{theorem}
\begin{proof}
Telescope $V^p_t(s):=\E_{a\sim p}Q^p_t(s,a)$ along the student's trajectory. For any $f$ with range at most
$u$, $|\E_qf-\E_pf|\le u\,\TV(p,q)$.
\end{proof}
\begin{remark}[What on-policy training buys]\label{rem:onpolicy}
\cref{thm:quadratic,thm:pdl} are the distillation versions of \citet{ross2010efficient} and
\citet{ross2011reduction}. On-policy training replaces $T^2\eoff$ by $uT\eon$. It removes compounding
only if (a)~the teacher is \emph{recoverable} ($u=O(1)$), and (b)~training actually makes $\eon$ small,
which DAgger-type no-regret arguments deliver relative to the best student \emph{on its own state
distribution}. In the non-recoverable cliff of \cref{thm:quadratic}, $u=T$, and no training distribution
avoids $\Theta(T^2\eps)$. For language models, $u$ is the worst-case cost of one bad token followed by
the teacher taking over. It is small for fluency-type costs and can be $\Theta(T)$ for commitments such
as a wrong early step in a proof. \Cref{sec:recoverability} makes this precise. $u$ is the best uniform constant,
it can be replaced by its average over the states where the student errs, and it is $O(1)$ when a
cost-sufficient latent chain of the teacher mixes, for example when the teacher repairs its own errors at a
constant rate. Related analyses of exposure bias include \citet{bengio2015scheduled},
\citet{ranzato2016sequence} and \citet{arora2022exposure}. On-policy distillation is studied by
\citet{lin2020autoregressive}, \citet{gu2024minillm}, \citet{agarwal2024gkd} and \citet{ko2024distillm}.
\end{remark}

\begin{theorem}[KL chain rules and compounding]\label{thm:kl-compounding}
\Pv
\begin{enumerate}[label=(\roman*),nosep]
\item Off-policy forward KL has $\cE_{\mathrm{shift}}=0$ for $\cE=\KL(P\|Q_\theta)$, and $\E\TV(P,Q_\theta)\le\sqrt{T\eoff^{\KL}/2}$.
\item For cumulative costs,
$|J_c(p)-J_c(q_\theta)|\le\sum_{t=1}^T\sqrt{\tfrac12\sum_{k\le t}e^{\mathrm{off},\KL}_k}\le T^{3/2}\sqrt{\eoff^{\KL}/2}$.
\item On-policy reverse KL has $\cE_{\mathrm{shift}}=0$ for $\cE=\KL(Q_\theta\|P)$. With recoverability,
$|J_c(q_\theta)-J_c(p)|\le uT\sqrt{\eon^{\mathrm{rKL}}/2}$.
\end{enumerate}
\end{theorem}
\begin{proof}
(i) and (iii) are \cref{thm:chain}. (ii) The law of $y_{\le t}$ has KL $\sum_{k\le t}e^{\mathrm{off},\KL}_k$ by
the truncated chain rule; apply Pinsker per $t$. The last claim is \cref{thm:pdl}.
\end{proof}
\noindent
Bound (ii) improves on \cref{thm:quadratic} only when errors are \emph{diffuse} (\cref{ex:regimes}(a)): there
$e^{\mathrm{off},\TV}\approx\sqrt{e^{\mathrm{off},\KL}}$, so $T^{3/2}\sqrt\eps$ beats $T^2\sqrt\eps$. For concentrated errors (deterministic
teacher, \cref{ex:regimes}(b)) the TV route is tighter. In the cliff with $T=100$ and $\eps=10^{-3}$, the true
gap is $4.9$, \cref{thm:quadratic} gives $5.05$, and (ii) gives $15.0$.

\begin{proposition}[Exposure-bias gap for mismatched pairings]\label{prop:exposure}
\Pv\Ck Let $g_t:\cS_t\to[0,G]$ be any bounded per-state loss, for example per-token TV ($G=1$),
$\JSD_\beta$ ($G=h(\beta)$), or clipped KL. Then
\[
\Big|\sum_{t=1}^T\big(\E_{\dstud t}-\E_{\dteach t}\big)g_t\Big|\le G\sum_{t=1}^T\TV(\dstud t,\dteach t)
\le G\sum_{t=1}^T\sum_{k<t}e^{\mathrm{off},\TV}_k\le G\,\tfrac{T(T-1)}2\max_ke^{\mathrm{off},\TV}_k.
\]
Also $\le G\sum_t\sqrt{\frac12\sum_{k<t}e^{\mathrm{off},\KL}_k}$. For a mixture prefix law
$\nu_t=\lambda\dstud t+(1-\lambda)\dteach t$, the gap is multiplied by $1-\lambda$. For \emph{unbounded} $g_t$,
for example reverse KL on teacher prefixes, no bound in terms of the training loss exists. The training
loss can be arbitrarily small while the on-policy loss is arbitrarily large, on states the teacher visits
rarely but the student visits often. With $T=2$, $p_1=(1-\eta,\eta)$, $q_1=(1-r\eta,r\eta)$, reverse KL $r^{3/2}$ at the
rare state and $\eta=r^{-2}$, the training loss is $O(r^{-1/2})$ while the on-policy loss is $\Theta(r^{1/2})$. Exactly
zero training loss does force $Q^x_\theta=\Pseq$ (\cref{prop:hybrid-div}).
\end{proposition}
\begin{proof}
Bounded-function duality for TV, then the hybrid argument for $\TV(\dstud t,\dteach t)$ as in \cref{prop:hybrid}.
\end{proof}

\begin{proposition}[Covariate-shift term]\label{prop:cov-term}
\Pv Let $\cE(\theta)=\E_x e_\theta(x)$ with $e_\theta\ge0$ \emph{pointwise}, for example a per-input divergence.
Under \cref{as:shift}, $\cE_{\mathrm{cov}}\le(\rho-1)\cE_\cD(\ttheta)$. If $0\le e_\theta\le G$, then
$\cE_{\mathrm{cov}}\le G\,\TV(\mu,\cD)$ (\cref{prop:covshift}). Pointwise non-negativity is needed: a task regret
can be negative at some inputs, and then the first bound fails.
\end{proposition}

\subsection{Recoverability: when is \texorpdfstring{$u=O(1)$}{u=O(1)}?}\label{sec:recoverability}

\Cref{thm:pdl} controls on-policy compounding through one worst-case number, $u$. This subsection shows that
$u$ is the right worst-case constant, replaces it by averages that can be measured, and gives conditions on
the teacher under which $u=O(1)$. The conditions are stated through the teacher's value functions and the
mixing of a Markov chain that it induces. Costs, $Q^p_t$ and $V^p_t$ are as in \cref{sec:shift}. The
\emph{local recoverability}
\[
u_t(s):=\max_aQ^p_t(s,a)-\min_aQ^p_t(s,a)=\max_aA^p_t(s,a)-\min_aA^p_t(s,a),\qquad A^p_t:=Q^p_t-V^p_t,
\]
is the range of the teacher's advantage at $s$, and $u=\max_t\sup_su_t(s)$. Let
$\TV_t(s):=\TV(p_t(\cdot\mid s),q_t(\cdot\mid s))$ and $\Delta_t(s):=\sum_a\big(q_t(a\mid s)-p_t(a\mid s)\big)Q^p_t(s,a)$. By
\cref{thm:pdl}, $\Delta J:=J_c(q_\theta)-J_c(p)=\sum_t\E_{\dstud t}\Delta_t$. Let $\nu^\theta$ (resp.\ $\nu^p$) be the law of $(t,s_t)$
when $t$ is uniform on $\{1,\dots,T\}$ and $s_t\sim\dstud t$ (resp.\ $s_t\sim\dteach t$). Then $\eon^{\TV}=\E_{\nu^\theta}\TV$.

\begin{theorem}[Recoverability where the student errs]\label{thm:local-u}
\Pv\Ck The following hold for every autoregressive student.
\begin{enumerate}[label=(\roman*),nosep]
\item \emph{Localisation.} $\Delta J=\bar a_{\mathrm{on}}\,T\eon^{\TV}$ and
$|\Delta J|\le\sum_t\E_{\dstud t}[u_t\TV_t]=\bar u_{\mathrm{on}}\,T\eon^{\TV}$, where, with $0/0:=0$,
\[
\bar a_{\mathrm{on}}:=\frac{\E_{\nu^\theta}\Delta}{\E_{\nu^\theta}\TV},\qquad
\bar u_{\mathrm{on}}:=\frac{\E_{\nu^\theta}[u\,\TV]}{\E_{\nu^\theta}\TV}\le u,\qquad
|\bar a_{\mathrm{on}}|\le\bar u_{\mathrm{on}}.
\]
So $\bar u_{\mathrm{on}}$ is the recoverability on the student's own prefixes, weighted by where the student errs.
\item \emph{Moments.} For $r\in[1,\infty]$ and $\frac1r+\frac1{r'}=1$,
\[
|\Delta J|\le T\,\|u\|_{L^r(\nu^\theta)}\,\|\TV\|_{L^{r'}(\nu^\theta)}\le T\,\|u\|_{L^r(\nu^\theta)}\,(\eon^{\TV})^{1-1/r}.
\]
\item \emph{Commitment states.} Fix $U\ge0$ and sets $B_t\supseteq\{s:u_t(s)>U\}$. Let
$\eps^B_{\mathrm{on}}:=\frac1T\sum_t\E_{\dstud t}[\TV_t\ind_{B_t}]$ be the on-policy error made on them, and
$\zeta:=\eps^B_{\mathrm{on}}/\eon^{\TV}$. Then
\[
|\Delta J|\le U\,T\eon^{\TV}+\sum_{t=1}^T(T-t+1)\,\E_{\dstud t}[\TV_t\ind_{B_t}]\le(U+\zeta T)\,T\eon^{\TV}.
\]
\item \emph{Sharpness.} Let $u_\cD:=\max_t\sup\{u_t(s):s=(x,y_{<t}),\ \cD(x)>0\}\le u$. \Cref{thm:pdl} holds with $u_\cD$ in
place of $u$. If every $p_t(\cdot\mid s)$ has full support, no smaller constant does:
$\sup|\Delta J|/(T\eon^{\TV})=u_\cD$ over autoregressive students with $\eon^{\TV}>0$. For each $(t,s)$ the ratio $u_t(s)$ is
attained by students with arbitrarily small $\eon^{\TV}$.
\end{enumerate}
\end{theorem}
\begin{proof}
(i) For any $f$ and laws $p,q$, $|\E_qf-\E_pf|\le\TV(p,q)(\max f-\min f)$, so $|\Delta_t|\le u_t\TV_t$. Sum over $t$, and
note that $\sum_t\E_{\dstud t}h_t=T\,\E_{\nu^\theta}h$.
(ii) Hölder's inequality, and $\TV^{r'}\le\TV$ because $\TV\le1$.
(iii) Split each expectation on $B_t$, and use $u_t\le T-t+1$, which holds because $Q^p_t\in[0,T-t+1]$.
(iv) See \cref{app:recoverability}.
\end{proof}

\noindent
By (iv), a bound that holds for every student cannot improve on $u$, even for students that are nearly
lossless. Any improvement must use \emph{where} the student errs, as (i)--(iii) do. Part (iii) makes
\cref{rem:onpolicy} quantitative. On-policy training is linear in $T$ when the student errs at states where
one token cannot cost much. It is quadratic when a fixed fraction $\zeta$ of its error falls on commitment
states. The average in (i) must be taken over the student's prefixes. An average over the teacher's
prefixes does not suffice:

\begin{proposition}[Teacher-side averages do not suffice]\label{prop:trap}
\Pv\Ck Let $\cV=\{0,1\}$, and let the prefix determine a latent state $z_t\in\{O,D,X\}$. Start at $z_1=O$. At $O$, token $0$
stays and token $1$ moves to $D$. At $D$, token $0$ returns to $O$ and token $1$ moves to $X$, which is absorbing. Step $t$
costs $\ind\{z_{t+1}\ne O\}$. The teacher always emits $0$. The student emits $1$ with probability $\eps\in(0,1]$ at $O$, always
at $D$, and never at $X$. Then $u_t=1$ at every prefix the teacher visits, so $\E_{\dteach t}u_t=1$ for all $t$ and
$\E_{\nu^p}[u\,\TV]/\E_{\nu^p}\TV=1$. But $u_t=T-t+1$ at $D$. Moreover $\eoff^{\TV}=\eps$ and
\[
\Delta J=\sum_{t=1}^T\big[1-(1-\eps)^t\big],\qquad T\eon^{\TV}=2-(1-\eps)^T-(1-\eps)^{T-1}.
\]
If $T\eps\le1$, then $\Delta J\ge\frac{T+1}8\,T\eon^{\TV}$, so $\bar u_{\mathrm{on}}\ge\frac{T+1}8$. By continuity, the same holds, up to factors
arbitrarily close to $1$, for a full-support teacher that emits $1$ with small enough probability.
\end{proposition}
\begin{proof}
See \cref{app:recoverability}.
\end{proof}

\noindent
The student never learned to leave $D$, a state the teacher never visits. This is exposure bias in its
purest form. Recoverability must therefore hold on the states that the student reaches. The next theorem
gives conditions on the teacher that guarantee this uniformly.

\begin{theorem}[Recoverability from forgetting]\label{thm:forgetting}
\Pv\Ck Each part relates $u$ to a property of the teacher alone.
\begin{enumerate}[label=(\roman*),nosep]
\item \emph{Value-function form.} For $k>t$ let $V^{(k)}_{t+1}(s'):=\E_p[c_k(s_k,y_k)\mid s_{t+1}=s']$ be the teacher's value
function for the single check $c_k$. Let $\iota_{t,k}(s):=\max_{a,a'}|V^{(k)}_{t+1}(s\,a)-V^{(k)}_{t+1}(s\,a')|$ be the largest
influence of the token emitted at $s$ on that value. Then
$u_t(s)\le\max_{a,a'}|c_t(s,a)-c_t(s,a')|+\sum_{k>t}\iota_{t,k}(s)$, and hence $u\le1+\max_t\sup_s\sum_{k>t}\iota_{t,k}(s)$.
\item \emph{Forgetting and coupling.} If $c_k(s_k,y_k)$ is a function of some statistic of $s_{k+1}$, then $\iota_{t,k}(s)$ is at
most the largest total variation between the teacher's laws of that statistic started from $s\,a$ and from
$s\,a'$. For any coupling of the teacher's continuations from $s\,a$ and $s\,a'$,
$|Q^p_t(s,a)-Q^p_t(s,a')|\le|c_t(s,a)-c_t(s,a')|+\E N$, where $N$ counts the steps $k>t$ at which the two
continuations incur different costs.
\item \emph{Mixing of a cost-sufficient chain.} Suppose there are maps $\phi_t:\cS_t\to\cZ_t$, with $\cZ_t$ countable, and $f_t$ with
$\phi_{t+1}(s\,a)=f_t(\phi_t(s),a)$, such that $p_t(\cdot\mid s)$ and $c_t(s,\cdot)$ depend on $s$ only through $z=\phi_t(s)$. The
teacher then induces a Markov chain on the latent states, with kernels
$K_t(z,z'):=\sum_ap_t(a\mid z)\ind\{f_t(z,a)=z'\}$. Let $\delta(K):=\sup_{z,z'}\TV(K(z,\cdot),K(z',\cdot))$ be Dobrushin's
contraction coefficient, and $\delta_m:=\max_{t\le T-m+1}\delta(K_t\cdots K_{t+m-1})$. For every $m$ with $\delta_m<1$,
\[
u\le1+\frac m{1-\delta_m}.
\]
If $c_t(s,a)$ depends on $(s,a)$ only through $\phi_{t+1}(s\,a)$, the cost of a step is a function of the latent state
it leads to, and then $u\le m/(1-\delta_m)$. Consequently $u\le1+2m_{1/2}$, where $m_{1/2}:=\min\{m:\delta_m\le\frac12\}$. For a time-homogeneous chain,
$m_{1/2}\le t_{\mathrm{mix}}(\frac14)$, its mixing time. Under a Doeblin condition,
$K_t\cdots K_{t+m-1}(z,\cdot)\ge\alpha\,\omega_t$ for all $t$ and $z$ and some probability laws $\omega_t$, one has $\delta_m\le1-\alpha$
and $u\le1+m/\alpha$.
\item \emph{The prefix chain itself never mixes.} For $\phi_t=\mathrm{id}$, $\delta(K_t\cdots K_{t+m-1})=1$ whenever $\cS_t$ has two
elements, because distinct prefixes have disjoint continuations. So $\delta_m=1$ for every $m<T$.
\end{enumerate}
\end{theorem}
\begin{proof}
(i) $Q^p_t(s,a)=c_t(s,a)+\sum_{k>t}V^{(k)}_{t+1}(s\,a)$, and the range of a sum is at most the sum of the ranges.
(ii) A function with values in $[0,1]$ changes its mean by at most the total variation. Under a coupling,
$Q^p_t(s,a)-Q^p_t(s,a')$ equals $c_t(s,a)-c_t(s,a')$ plus the expected sum of the later cost differences. Each
difference lies in $[-1,1]$ and vanishes unless $N$ counts it.
(iii) See \cref{app:recoverability}.
(iv) $K_t\cdots K_{t+m-1}(s,\cdot)$ is supported on extensions of $s$.
\end{proof}

\noindent
Part (i) says that $u$ is small when no single token moves the teacher's probability of failing a later
check by much, summed over the later checks. Part (ii) makes this measurable. Force $a$ and $a'$ at $s$, roll
the teacher out from both with shared sampling noise, and count the steps whose costs differ. Part (iii)
turns the condition into mixing. The chain must live on a \emph{cost-sufficient} summary $z$ of the prefix, for
example the state of a format or unit-test checker together with whatever the teacher's next-token law
depends on. By (iv), the prefix tree itself never mixes. Doeblin's condition says the following. From every
latent state, including erroneous ones, the teacher reaches within $m$ tokens, with probability at least $\alpha$,
a state whose law does not depend on where it started. This is \emph{self-correction}. Then $u-1$ is at most
$m/\alpha$, which bounds the expected time that self-correction takes.

\begin{example}[Commitment versus self-correction]\label{ex:selfcorrect}
\Pv\Ck Take two latent states, $C$ (on track) and $E$ (in error), and let each step that ends in $E$ cost $1$. From $C$ the
teacher moves to $E$ with probability $\eta$. From $E$ it repairs, moving to $C$, with probability $\alpha$. Let
$\alpha+\eta\le1$, and let some token lead from each state to each state. Then $\delta_1=\lambda_2:=1-\alpha-\eta$, the second
eigenvalue of $K$, and
\[
u=\frac{1-\lambda_2^T}{1-\lambda_2}\ \nearrow\ \frac1{\alpha+\eta}\qquad(T\to\infty).
\]
So the bound of \cref{thm:forgetting}(iii) for costs that depend only on the next latent state is attained in
the limit. Here $u$ is the relaxation time of the teacher's error process, close to the expected repair time
$1/\alpha$ when $\eta\ll\alpha$. Without repair and without teacher errors ($\alpha=\eta=0$), this is the cliff of
\cref{thm:quadratic}: $\delta_m=1$ for all $m$, and $u=T$.
\end{example}
\begin{proof}[Computation]
Let $\bar W_j(z)$ be the expected cost of the step that ends in $z$ at time $j$, plus that of all later steps. Then
$u_t(s)=D_{t+1}:=\bar W_{t+1}(E)-\bar W_{t+1}(C)$, with $D_{T+1}=1$ and $D_{t+1}=1+\lambda_2D_{t+2}$. So
$D_{t+1}=\sum_{j=0}^{T-t}\lambda_2^j$, which is largest at $t=1$.
\end{proof}

\begin{proposition}[The value martingale, and a teacher-side bound]\label{prop:valmart}
\Pv\Ck Let $C:=\sum_{t=1}^Tc_t(s_t,y_t)$, $\Var_p(C):=\E_{x\sim\cD}\Var_{y\sim\Pseq}(C)$, and
$\sigma_t^2(s):=\Var_{a\sim p_t(\cdot\mid s)}Q^p_t(s,a)\le u_t(s)^2/4$.
\begin{enumerate}[label=(\roman*),nosep]
\item Under the teacher, $G_t:=\sum_{k<t}c_k(s_k,y_k)+V^p_t(s_t)$ is the Doob martingale of $C$. Its increments are
$A^p_t(s_t,y_t)$, so
\[
\Var_p(C)=\sum_{t=1}^T\E_{\dteach t}\sigma_t^2=T\,\E_{\nu^p}\sigma^2 .
\]
\item For a terminal reward $r$ (costs $0$ before step $T$, and $c_T=1-r$), $Q^p_t(s,a)=1-\E_p[r\mid s_{t+1}=s\,a]$. So $u_t(s)$
is the largest one-token swing of the teacher's expected reward, which is its success probability for
$0/1$ rewards. Hence $u\le1$ and $|J_r(p)-J_r(q_\theta)|\le\sum_t\E_{\dstud t}[u_t\TV_t]\le T\eon^{\TV}$. Also
$\Var_p(C)\le\E_x[\bar r(x)(1-\bar r(x))]\le\frac14$, where $\bar r(x):=\E_p[r\mid x]$.
\item $|\Delta J|\le\sqrt{\E_x\chi^2(\Qseq\|\Pseq)\,\Var_p(C)}$. If $\chi^2(q_t(\cdot\mid s)\|p_t(\cdot\mid s))\le\chi^2_\infty$ for all $t$ and $s$, then
$\E_x\chi^2(\Qseq\|\Pseq)\le(1+\chi^2_\infty)^T-1$. If moreover $T\chi^2_\infty\le1$, then
\[
|\Delta J|\le T\sqrt{2\chi^2_\infty\,\E_{\nu^p}\sigma^2}.
\]
\end{enumerate}
\end{proposition}
\begin{proof}
(i) $G_{t+1}-G_t=c_t(s_t,y_t)+V^p_{t+1}(s_{t+1})-V^p_t(s_t)=A^p_t(s_t,y_t)$. Given $s_t$, it has mean $0$ and variance
$\sigma_t^2(s_t)$. Also $G_1=\E_p[C\mid x]$ and $G_{T+1}=C$, and martingale increments are orthogonal. The bound
$\sigma_t\le u_t/2$ is Popoviciu's inequality.
(ii) The formula for $Q^p_t$ is the definition; then apply \cref{thm:local-u}(i). Moreover
$\Var(r\mid x)\le\bar r(1-\bar r)$ because $r^2\le r$.
(iii) For each $x$, $\E_{\Qseq}C-\E_{\Pseq}C=\E_{\Pseq}\big[(\Qseq/\Pseq-1)(C-\E_{\Pseq}C)\big]$. Apply Cauchy--Schwarz, and then again over $x$.
Summing out the last token gives $\sum_y\Qseq(y)^2/\Pseq(y)\le(1+\chi^2_\infty)\sum_{y_{<T}}\Qseq(y_{<T})^2/\Pseq(y_{<T})$; induct on $T$.
Finally, $(1+\chi^2_\infty)^T-1\le e^{T\chi^2_\infty}-1\le2T\chi^2_\infty$ when $T\chi^2_\infty\le1$, and $\Var_p(C)=T\,\E_{\nu^p}\sigma^2$ by (i).
\end{proof}

\begin{remark}[Where the average is taken]\label{rem:where-avg}
The results above settle when $u$ can be replaced by an average.
\begin{center}\small
\begin{tabular}{@{}>{\raggedright\arraybackslash}p{0.27\linewidth}>{\raggedright\arraybackslash}p{0.22\linewidth}>{\raggedright\arraybackslash}p{0.43\linewidth}@{}}
\toprule
recoverability measured & student's error & bound on $|\Delta J|$\\\midrule
worst case, $u$ & on-policy mean, $\eon^{\TV}$ & $uT\eon^{\TV}$; sharp (\cref{thm:local-u}(iv))\\
student's prefixes, $\bar u_{\mathrm{on}}$ & on-policy mean & $\bar u_{\mathrm{on}}T\eon^{\TV}$ (\cref{thm:local-u}(i))\\
teacher's prefixes, $\E_{\nu^p}\sigma^2$ & uniform, $\chi^2_\infty$ & $T\sqrt{2\chi^2_\infty\E_{\nu^p}\sigma^2}$ if $T\chi^2_\infty\le1$ (\cref{prop:valmart})\\
teacher's prefixes & on-policy mean & none (\cref{prop:trap})\\\bottomrule
\end{tabular}
\end{center}
A teacher-side average suffices when the student is accurate everywhere, including at states the teacher
avoids (\cref{prop:valmart}). Accuracy on the teacher's own prefixes is not enough (\cref{prop:trap}). The
teacher-side quantity is $\Var_p(C)$, the variance of the teacher's own total cost, which can be estimated by
sampling from the teacher alone. $\Var_p(C)=O(T)$ is a weak-dependence condition on
the teacher's cost process. It holds whenever $u=O(1)$, because $\Var_p(C)\le Tu^2/4$. By contrast, an absorbing
error that the teacher enters with probability of order $1/T$ per step makes $\Var_p(C)$ of order $T^2$.
\end{remark}

\begin{remark}[Turning the bound into a prediction]\label{rem:measure-u}
$\bar u_{\mathrm{on}}$ and $\bar a_{\mathrm{on}}$ can be estimated from the two models and teacher roll-outs. Sample $y\sim\Qseq$, compute $\TV_t(s_t)$
exactly at every position, and draw one position $t$ with probability proportional to $\TV_t(s_t)$. Let
$S:=\sum_{t'}\TV_{t'}(s_{t'})$. Then $S\,u_t(s_t)$ is unbiased for $\bar u_{\mathrm{on}}T\eon^{\TV}$, and $S\,\Delta_t(s_t)/\TV_t(s_t)$ is unbiased for
$\Delta J$. At the drawn $s_t$, estimate $Q^p_t(s_t,a)$ for the $k$ tokens that carry most of the mass of $|q_t-p_t|$, from
$n$ teacher continuations each. By Hoeffding's inequality and a union bound, all $k$ estimates are within
$(T-t+1)\sqrt{\log(2k/\delta)/(2n)}$ with probability at least $1-\delta$. Let $m_t$ be the mass of $|q_t-p_t|$ on the
other tokens. Then $|\Delta_t|$ is at most $\TV_t$ times the range of $Q^p_t(s_t,\cdot)$ over the $k$ tokens, plus
$(T-t+1)\,m_t$. So that range can replace $u_t(s_t)$ at the cost of the residual, and the same roll-outs
estimate $\Delta_t(s_t)$ up to the Monte Carlo error plus $(T-t+1)\,m_t$. The
prediction is $\Delta J\approx\bar a_{\mathrm{on}}T\eon^{\TV}$. Its growth with $T$ follows $\bar u_{\mathrm{on}}$: it is linear while $\bar u_{\mathrm{on}}$
stays bounded, and quadratic when the commitment fraction $\zeta$ of \cref{thm:local-u}(iii) stays bounded away
from $0$.
\end{remark}


\subsection{Assembling the bound, and which term dominates}

\begin{corollary}[End-to-end bound for off-policy soft-label forward-KL distillation]\label{cor:e2e-offpolicy}
\Pa{\cref{as:bounded,as:shift}} Let $\cA=\min_\theta\E_\cD\KL(P\|Q_\theta)$. With probability $\ge1-\delta$,
\[
\E_\mu\TV(\Pseq,Q^x_{\ttheta})\le\sqrt{\frac\rho2\Big[\cA+2\Big(2\Rad_n+TB\sqrt{\tfrac{\log(2/\delta)}{2n}}\Big)+\cE_{\mathrm{opt}}\Big]}.
\]
Under \cref{as:real,as:finite}, with hard labels and exact ERM, instead
$\E_\mu\TV\le\min\big\{\rho\,\eta_n,\ \eta_n+\TV(\mu,\cD)\big\}$ with $\eta_n=\sqrt{2(R\ln2+\log(1/\delta))/n}$.
\end{corollary}
\begin{proof}
\cref{thm:master} with $\cE_{\mathrm{shift}}=0$ (\cref{thm:chain}), \cref{thm:unif}, \cref{prop:cov-term},
and Pinsker. For the second claim use \cref{thm:mle} and \cref{prop:covshift}.
\end{proof}

\begin{heuristic}[The dominant term in large-scale LLM distillation]\label{heur:dominant}
(1)~$\cE_{\mathrm{est}}$ is small. Teacher data is effectively unlimited, soft labels remove label noise
(\cref{prop:rb,prop:exact}), and realisable rates are horizon-free (\cref{thm:mle}).
(2)~$\cE_{\mathrm{opt}}$ is real but shrinks with compute (\cref{thm:sgd}); it dominates only for
undertrained students.
(3)~$\cA$ is the floor. It dominates for small students (\cref{thm:bottleneck}; \cref{sec:lower}) and, at
sequence level, grows linearly in $T$ (\cref{cor:approx-T}).
(4)~$\cE_{\mathrm{shift}}$ is zero for the exact pairings, but for mismatched pairings and cumulative
criteria it \emph{amplifies} $\cA$ by up to $O(T)$ (\cref{thm:quadratic,prop:exposure}). On-policy training removes
this amplification only for the part of the student's error made at recoverable states (\cref{thm:local-u}).
Hence in the capacity-limited regime the dominant term is \emph{approximation error as seen by the
deployment criterion}: the irreducible floor, amplified by prefix shift and \emph{misallocated} by the
choice of divergence (\cref{prop:agnostic}). \Cref{sec:objectives} designs an objective against exactly
this term.
\end{heuristic}

